\begin{exercise}[Exercise 10-10 from \cite{lee2011smooth}]
Let $M$ be a manifold without boundary of dimension $n$ and $E \to M$ a rank $k$ vector bundle over $M$. Then there exists a section $\sigma \colon M \to E$ with the following properties:
\begin{enumerate}
    \item If $k > n$, then $\sigma$ is nowhere vanishing
    \item If $k \leq n$, then the zeros of $\sigma$ form a closed codimension $k$ submanifold of $M$
\end{enumerate}
In particular, there exists a smooth vector field $X$ on $M$, where the zeros of $X$ are a discrete closed subset. Hence, every compact manifold with or without boundary admits a smooth vector field with only finitely many zeros. 
\end{exercise}

\begin{proof}
    Consider the zero section $\xi \colon M \to E$ and let $M^\prime := \mathrm{im}(\xi)$. Then $M^\prime$ is a codimension $0$ submanifold (without boundary) of $E$. By the transversality homotopy theorem (Theorem \ref{thm:transversality_homotopy}) any section $M \to E$ is homotopic to a section $\sigma \colon M \to E$ that is transversal to $M^\prime$. If $k > n$ this is equivalent to 
    \[
    \sigma(M) \cap M^\prime = \emptyset,
    \]
    which in turn is equivalent to $\sigma(p) \neq 0$ for all $p \in M$. If $k \leq M$, then 
    \[\{p \in M \mid \sigma(p) = 0\} = \sigma^{-1}(M^\prime) \]
    and $\sigma^{-1}(M^\prime)$ is a closed codimension $k$ submanifold of $M$. This shows the first two statements. Applying these results to $E := TM$, we see that the zeros of $\sigma$ are a closed codimension $0$ submanifold, thus discrete. If $M$ is compact this is equivalent to $\{p \in M \mid \sigma(p) = 0\}$ being finite. If $M$ has non-empty boundary, we can construct a vector field with finitely many zeros on the double and restrict it to $M$.  
\end{proof}